\section{Threats to Validity}
\label{sec:threats}

\paragraph{Internal threats.} Every arm runs 5 trials with the same per-trial
seeds on the same build, schema, executor, oracles and budget, each trial has
its own server, a metamorphic mismatch counts only if it repeats, and every
candidate is reproduced three times before triage. But five trials are few:
differences with $p$ above 0.008 are effect sizes, not significant results, and
several conclusions rest on them.

Our own oracles can be wrong. Ten false-positive classes reached triage, six
before the campaigns and four in them (Section~\ref{sec:rq2}). We fixed the
last four \emph{after} the campaigns and checked them on the stored
reproducers and with unit tests, but did not re-run the campaigns, so the
precision after the fixes is not a measurement. Six candidates are open and more
classes may remain. Verdicts are ours, some made from the reduced reproducer; we
release the reproducers. An accepted LLM rewrite is tested, not proven
(Section~\ref{sec:rq4}). Attribution in \rqref{3} replays on single-bug servers,
so a finding that needs two injected bugs is reported as unattributed.

The two laptops differ, so \rqref{1} and \rqref{3} are not comparable in
absolute throughput. One worker's server failed crash recovery during \rqref{1}
and 14 trials were re-run with the same seeds; a laptop suspension leaves a gap
in wall-clock time but not in the monotonic one-hour budget.

\paragraph{External threats.} We test one system at one commit, with one schema
of 9 relations and one set of row-level security policies that we chose;
bugs that need other data shapes, extensions or concurrent sessions are out of
reach. The 23 benchmark bugs are recent fixes with regression tests and may
favour features that recent \pg development touched. The LLM results hold for
free-tier Gemini Flash models on 10 commits; we store every prompt and response.

\paragraph{Construct threats.} Coverage measures what code ran, not whether an
oracle could have seen a bug there. A bug ``found'' in \rqref{3} means a test
case failed on its solo server and passed on the fixed one, even if its symptom
differs from the PoC's. A fragment that passes the \gate reaches the changed
lines; that does not prove it tests the changed behaviour.
